\section{Architectural Evolution toward Selective Representation Sharing}

% \subsection{Cross-Stitch Networks}
% Cross-Stitch shifted the degree of representation sharing from a manually
% specified design choice to a learnable component of the model. As shown in
% Figure~\ref{fig:cross-stitch}, given task-specific representations $x_A$ and
% $x_B$, a cross-stitch unit learns the degree of representation sharing through
% linear combinations $\tilde{x}_A$ and $\tilde{x}_B$:
% \begin{equation}
% \begin{bmatrix}
% \tilde{x}_A \\
% \tilde{x}_B
% \end{bmatrix}
% =
% \begin{bmatrix}
% \alpha_{AA} & \alpha_{AB} \\
% \alpha_{BA} & \alpha_{BB}
% \end{bmatrix}
% \begin{bmatrix}
% x_A \\
% x_B
% \end{bmatrix},
% \label{eq:cross-stitch}
% \end{equation}
% where the resulting representations are fed into subsequent task-specific
% layers \citep{Misra2016CrossStitch}.

% Despite its limited sample adaptivity due to static mixing weights shared
% across samples \citep{Tang2020PLE}, and the additional parameter cost
% associated with maintaining task-specific network streams
% \citep{Ma2018MMoE}, Cross-Stitch provided an architectural insight relevant
% to multi-objective recommendation: representation sharing need not be
% architecturally fixed or binary, but can instead be learned while preserving
% task-specific representations.


% \begin{figure}[t]
%     \centering
%     \fbox{
%         \parbox[c][4.5cm][c]{0.85\linewidth}{
%             \centering
%             Placeholder for Figure 2\\[0.5em]
%             Cross-Stitch Network\\[0.5em]
%             Two task-specific streams with learnable\\
%             cross-task representation mixing
%         }
%     }
%     \caption{Conceptual illustration of the Cross-Stitch architecture, where task-specific representations are selectively mixed through learnable cross-stitch units.}
%     \label{fig:cross-stitch}
% \end{figure}


\subsection{Multi-gate Mixture-of-Experts (MMoE)}

Building on the shared-bottom model, MMoE has become a representative architecture for industrial multi-task recommendation \citep{Ma2018MMoE}. Figure~\ref{fig:architectures}(b) illustrates the architecture, which retains a task-specific tower for each task, with the prediction for task $k$ given by

\begin{equation}
y_k = h^k\left(f^k(x)\right),
\label{eq}
\end{equation}

where $h^k(\cdot)$ denotes the task-specific tower and $f^k(x)$ is the representation supplied to it. Rather than producing this representation through a single shared-bottom network, MMoE combines a pool of shared experts:

\begin{equation}
f^k(x)
=
\sum_{i=1}^{n} g_i^k(x)f_i(x),
\label{eq}
\end{equation}

where $f_i(x)$ is the output of expert $i$ and $g_i^k(x)$ is its weight for task $k$. These weights are generated by an input-dependent, task-specific gate:

\begin{equation}
g^k(x)
=
\operatorname{softmax}\left(W_{gk}x\right).
\label{eq}
\end{equation}

MMoE therefore introduces task-specific selection at the gate level, allowing each task to form input-dependent expert mixtures while keeping all expert parameters shared. In the original experiments, MMoE showed its clearest advantage over baselines when tasks were weakly related and also improved performance in large-scale content recommendation \citep{Ma2018MMoE}. However, insufficient differentiation among shared experts may still produce the seesaw phenomenon, in which improvement in one task accompanies degradation in another \citep{Tang2020PLE}. This limitation motivates PLE to introduce task-specific experts.

\subsection{Progressive Layered Extraction (PLE)}
To address the limited expert differentiation in MMoE, PLE explicitly separates shared and task-specific experts \citep{Tang2020PLE}. As Figure~\ref{fig:architectures}(c) shows, a task-specific gate at each extraction layer combines its corresponding task-specific experts with the shared experts, while a shared gate aggregates information from all experts for the next layer. Stacking these extraction modules progressively separates shared and task-specific representations. For task $k$, the output of its gate at extraction layer $j$ is

\begin{equation}
g^{k,j}(x)
=
w^{k,j}\left(g^{k,j-1}(x)\right)S^{k,j}(x),
\label{eq}
\end{equation}

where $w^{k,j}(\cdot)$ produces gate weights from the previous-layer representation and $S^{k,j}(x)$ contains the expert outputs available to task $k$ at layer $j$. Experiments on large-scale video recommendation data show that PLE outperforms MMoE across different task relationships and task-group sizes \citep{Tang2020PLE}. However, PLE retains shared embeddings, which may obscure task-specific user preferences \citep{Su2024STEM}. Moreover, adding a task requires additional task-specific experts, gates, and a tower, causing task-specific capacity to grow with the number of tasks. 


% \subsection{Shared and Task-specific Embeddings (STEM)}
% To address this limitation, STEM incorporates shared and
% task-specific embeddings to capture both common and
% task-specific user preferences \citep{Su2024STEM}.
% As illustrated in Figure~\ref{fig:stem}, STEM-Net allows
% each task tower to receive outputs from shared and all
% task-specific experts, while blocking gradients through
% other tasks' expert outputs. The representation for task
% $t$ is computed as
% \begin{equation}
% \begin{aligned}
% o^t ={}&
% \sum_{i=1}^{K_1} g_i^{t\to t}h_i^t
% +
% \sum_{j=1}^{K_2} g_j^{S\to t}h_j^S \\
% &+
% \sum_{\substack{t'=1 \\ t'\neq t}}^{T}
% \sum_{i=1}^{K_1}
% g_i^{t'\to t}\operatorname{SG}(h_i^{t'}),
% \end{aligned}
% \label{eq:stem-gating}
% \end{equation}
% where $h$ denotes expert outputs and $g$ denotes gating
% weights. Each of the $T$ tasks has $K_1$ task-specific
% experts, alongside $K_2$ shared experts.
% The stop-gradient operator $\operatorname{SG}$ preserves
% outputs in the forward pass but blocks gradients in
% backpropagation, enabling cross-task knowledge transfer
% while protecting task-specific experts and embeddings
% from updates by other tasks.

% In their experiments, STEM-Net outperforms PLE in average AUC across all three datasets and achieves positive transfer relative to single-task models on the Like task in both TikTok and QK-Video, supporting the effectiveness of its selective-sharing design.


\subsection{Shared and Task-specific Embeddings (STEM)}
To address the shared-embedding limitation retained by PLE, STEM introduces one shared embedding table and a separate embedding table for each task, enabling common and task-specific user preferences to be represented separately \citep{Su2024STEM}. Shared experts receive the shared embeddings, whereas each task-specific expert group receives embeddings from its corresponding task-specific table. As illustrated in Figure~\ref{fig:architectures}(d), each task tower then aggregates outputs from its own task-specific experts, the shared experts, and the other tasks' experts, weighted by gates computed from both shared and task-specific embeddings. Outputs from the other tasks' experts are passed through a stop-gradient operator, which preserves them in the forward pass but blocks their gradients, enabling direct cross-task transfer while protecting task-specific parameters from updates by other tasks.

STEM-Net outperforms PLE in average AUC across all three datasets and achieves positive transfer over the corresponding single-task baselines on the low-positive-rate Like task in both TikTok and QK-Video \citep{Su2024STEM}. However, assigning task-specific embeddings to additional feature fields can substantially increase parameter count, motivating the capacity-allocation question examined later.

% As illustrated in Figure~\ref{fig:architectures}(d), each task tower then receives outputs from the shared and all task-specific experts. The representation for task $t$ is

% \begin{equation}
% \begin{aligned}
% o^t ={}&
% \sum_{i=1}^{K_1} g_i^{t\to t}h_i^t
% +
% \sum_{j=1}^{K_2} g_j^{S\to t}h_j^S \
% &+
% \sum_{\substack{t'=1 \ t'\neq t}}^{T}
% \sum_{i=1}^{K_1}
% g_i^{t'\to t}\operatorname{SG}(h_i^{t'}),
% \end{aligned}
% \label{eq:stem-gating}
% \end{equation}

% where $h$ denotes expert outputs and $g$ denotes gating weights computed from both shared and task-specific embeddings. Each of the $T$ tasks has $K_1$ task-specific experts, alongside $K_2$ shared experts. The stop-gradient operator SG preserves other tasks' expert outputs in the forward pass but blocks their gradients, enabling direct cross-task transfer while protecting task-specific parameters from updates by other tasks.

\subsection{Recent Diversification of Selective Sharing}
The progression above assumes that task-specific capacity, once added, can be effectively trained and served at acceptable cost.
HoME \citep{wang2025home} identifies expert collapse, degradation, and underfitting,
showing that data-sparse tasks—precisely those that motivate task-specific capacity—may
fail to utilise their own experts. MoME \citep{Xu2024MoME} and SMES \citep{Zhang2026SMES}
instead question its cost, reducing storage redundancy through masked expert
subnetworks and limiting the experts executed at inference through sparse routing.
Because these studies optimise different efficiency measures, existing results do not
establish which placement of task-specific capacity best preserves task-level positive
transfer under comparable parameter budgets.




